\documentclass[../../../include/open-logic-section]{subfiles}
%READERNOTE{SOL6-A245: मूळाच्या प्रतिमा-तुलनेत निवडीची पुरेशी अट स्पष्ट केली. आच्छादक फलनापासून उलट दिशेचे एकास-एक फलन निवडीच्या स्वयंसिद्धकाखाली मिळते; प्रतिस्थापनाला स्वतः निवड आवश्यक आहे किंवा ही तुलना पूर्ण निवडीशी समतुल्य आहे असा दावा नाही.}

\begin{document}

\olfileid{sth}{replacement}{limofsize}
\olsection[आकारमानाची मर्यादा]{आकारमानाच्या मर्यादेचे तत्त्व}

प्रतिस्थापनाचे ``आंतरिक'' समर्थन देण्याचा बहुधा सर्वाधिक प्रचलित प्रयत्न पुढील कल्पनेवर आधारलेला असतो:
\begin{enumerate}
  \item[] \limofsize. कोणत्याही वस्तूंचा संच तयार होतो, जर त्या फार जास्त नसतील तर.
\end{enumerate}
आकारमानाच्या पुढील तुलनेसाठी निवडीचे स्वयंसिद्धक गृहीत धरून, हे तत्त्व प्रतिस्थापनाला लगेच समर्थन देईल. कारण प्रतिस्थापनाने तयार झालेला कोणताही संच ज्या संचापासून तयार झाला त्याच्यापेक्षा मोठा असू शकत नाही. नेमके सांगायचे तर: समजा प्रतिस्थापन वापरून तुम्ही
$\funimage{\tau}{A} = \Setabs{\tau(x)}{x \in A}$ हा संच तयार केला; मग $\cardle{\funimage{\tau}{A}}{A}$; म्हणून $A$ चे !!{element}s संच तयार होण्याइतके कमी असतील, तर त्यांच्या प्रतिमाही $\funimage{\tau}{A}$ तयार होण्याइतक्या कमी असतील.

प्रतिस्थापनाचे समर्थन करण्यासाठी \limofsize{} वापरण्यातील उघड अडचण अशी की \limofsize{} सारखे कोणतेही तत्त्व आपण अजून मांडलेलेच \emph{नाही}. शिवाय, \crefrange{sth:story::chap}{sth:z::chap} मध्ये संचांच्या संचयी-क्रमिक संकल्पनेची मांडणी करताना \limofsize{} कडे निर्देश करणारा \emph{संकेतही} त्यात नव्हता. म्हणूनच बूलोस यांनी एकदा असे लिहिले: ``संच उपपत्तीच्या `मागे' निदान दोन विचार आहेत, असा निष्कर्ष कदाचित काढता येईल'' \citeyearpar[p.~19]{Boolos1989}. एका बाजूला, संचांच्या संचयी-क्रमिक संकल्पनेशी संबंधित कल्पना $\Z$ ला समर्थन देण्यासाठी आहेत. दुसऱ्या बाजूला, \limofsize{} हे प्रतिस्थापनाला समर्थन देण्यासाठी आहे.

पण अडचण एवढीच नाही की \limofsize{} बाबत आपण आतापर्यंत \emph{मौन} बाळगले आहे. उलट, नुकतेच मांडलेले \limofsize{} तत्त्व संचांच्या संचयी-क्रमिक संकल्पनेशी जुळणेही अवघड वाटते. कारण संच \emph{टप्प्यांत} तयार होतात, या कल्पनेचा त्यात अजिबात उल्लेख नाही.

या ``दोन विचारांचा'' (म्हणजे संचांची संचयी-क्रमिक संकल्पना आणि \limofsize) इतिहास पाहता हे फारसे आश्चर्यकारक नाही. संच उपपत्तीतील विरोधाभास रोखण्याचे हे मुळात दोन भिन्न प्रकल्प आहेत. संचांची संचयी-क्रमिक संकल्पना रसेलचा विरोधाभास साधारणपणे असे म्हणून रोखते: \emph{रसेलचा संच अस्तित्वात असेल अशी अपेक्षाच आपण करायला नको होती, कारण तो कोणत्याही टप्प्यावर ``तयार'' होणार नाही}. याउलट, \limofsize{} रसेलचा संच साधारणपणे असे म्हणून नाकारते: \emph{रसेलचा संच अस्तित्वात असेल अशी अपेक्षाच आपण करायला नको होती, कारण तो फार मोठा असता}.

असे मांडल्यावर स्पष्टच सांगू: विरोधाभासांना उत्तर म्हणून \limofsize{} ला \emph{खूपच} अधिक समर्थन आवश्यक आहे. उदाहरणार्थ, रसेलच्या विरोधाभासाचे हे रूप पाहा: \emph{स्वतःचाच वास न घेणाऱ्या पग कुत्र्यांचाच नेमका वास घेणारा कोणताही पग कुत्रा नाही} (\olref[sth][story][rus]{sec} पाहा). ``असा पग कुत्रा का नाही?'' असे विचारल्यावर त्याला फारच जास्त पग कुत्र्यांचा वास घ्यावा लागला असता, हे चांगले उत्तर ठरणार नाही. मग रसेलचा संच अस्तित्वात नाही याचे सहज पटणारे स्पष्टीकरण तो अस्तित्वात असण्यासाठी ``फार मोठा'' झाला असता, हे कसे ठरेल?

थोडक्यात, \limofsize{} मागील ``सहजस्फूर्त'' प्रेरणेविषयी तुम्हाला थोडे गूढ वाटले, तर ते समजण्यासारखे आहे.

\end{document}
